\section{Why This Graph Got Away With It}
\label{sec:ch11-why-this-graph-got-away-with-it}

\index{representation order!the ascending coincidence|(}%
The pages that still answer correctly are the conference schedule's doing, and
nothing else.
\code{sessions} comes back in start-time order, and the four sessions in that
order belong to speakers one, one, two and three.
Chapter~\ref{ch:router-n-plus-1} established that the router de-duplicates
before it sends, so the representation list is speakers one, two and three.
SQLite, asked for exactly those three by primary key, hands the rows back
ascending: one, two, three. The list the subgraph was asked about and the list
the store returned are the same list, so pairing them by position happens to
be right.

\index{representation order!a property of the data, not of the code}%
That is a property of the seed data. It is not a property of the code, nobody
chose it, and it survives no change at all: move one talk to a slot that puts
its speaker out of turn and the answers start disagreeing with the questions.

\index{representation order!no page in this graph can catch it}%
It also means no paging request in this graph can find the defect. The parent
list has to arrive in an order the subgraph's store will not reproduce, and
\code{sessions} sorts by \code{StartsAt} inside the resolver and takes no sort
argument, so the only orders a client can ask for are windows of that one
order. Relay's \code{last:} does not help either: it selects the final entries
and still hands them back forwards. The request that finds it is the one where
the client supplies the list, which in this schema is \code{nodes(ids:)}, and
that is why the chapter opened on a field a smoke test rarely touches.

\index{representation order!aliased node calls are immune}%
Two aliased \code{node(id:)} calls, which look like the same request written
differently, are immune. The planner treats them as two separate entity
fetches, each carrying one representation, and a batch of one cannot be out of
order. Both come back correct on this graph. So the failing request and the
passing request differ only in whether the client wrote one root field or two.

\subsection{Your Graph Is Probably Less Lucky}

\index{representation order!when the orders will not agree}%
Two shapes make the coincidence go away, and between them they cover most of
what a federated graph is for.

\index{representation order!a parent list sorted by something else}%
The parent list is sorted by something that is not the child's key. A feed by
date, search results by relevance, a leaderboard by score: every one of them
produces entity keys in an order the child's own store has never heard of. This
is the normal case, and this book's schedule is only outside it because the
speakers were seeded in the order they speak.

\index{representation order!a store with no natural order}%
The child store has no order to be ascending in. A key-value cache, a document
database, a search index, a service behind HTTP that answers a list of ids: none
of these promises anything about sequence, and several of them will return the
cheapest order rather than a stable one. The moment a loader talks to one of
those, the code that looked fine against SQLite is answering a coin toss.

\subsection{One Half Is Loud and the Other Is Silent}

\index{DataLoader!breaking the first clause}%
The loader in section~\ref{sec:ch11-two-sessions-each-others-speakers} breaks
both of \code{FetchAsync}'s clauses, and the graph reports one of them.

Send it three representations where the middle key matches nothing, on the
service's own port so that no router is involved, and ask for the key back
alongside the name:

\begin{minted}[fontsize=\footnotesize,breakanywhere]{graphql}
query E($r: [_Any!]!) {
  _entities(representations: $r) { ... on Speaker { id name } }
}
\end{minted}

\begin{minted}[fontsize=\footnotesize,breakanywhere]{json}
{"r":[{"__typename":"Speaker","id":"U3BlYWtlcjox"},
      {"__typename":"Speaker","id":"U3BlYWtlcjo5OQ=="},
      {"__typename":"Speaker","id":"U3BlYWtlcjoy"}]}
\end{minted}

\begin{minted}[fontsize=\footnotesize,breakanywhere]{json}
{"data":{"_entities":[
   {"id":"U3BlYWtlcjox","name":"Ada Fischer"},
   {"id":"U3BlYWtlcjoy","name":"Bruno Kaminski"},
   null]}}
\end{minted}

\index{DataLoader!an unwritten slot is an error}%
That is the data half. Beside it the response carries an \code{errors} array,
and the message in it is the one this book keeps meeting:
\code{Unexpected Execution Error}, on the entity route. Whether the path also
names the position that was never filled is not something to build on. Most
runs give the index and some do not, and I have no account of why.

\index{DataLoader!an unresolved result is not a null}%
Two rows came back for three keys, the loop ran twice, and
\code{results.Span[2]} was never written at all. An unwritten slot is not a
null: it is a \code{Result} that was never resolved, and the engine turns it
into an execution error rather than into the null the specification asks for.
Which is the fourth distinct cause of \code{Unexpected Execution Error} in this
book, after a stub advertising a key it cannot resolve, a service missing its
node id serializer, and the dictionary with a duplicate key in it that
chapter~\ref{ch:router-n-plus-1} ended on. The message still never says which.

\index{DataLoader!the loud half is the survivable one}%
Loud is a relative word here. The service writes nothing while this happens:
no exception, no stack trace, nothing on standard output or standard error but
the statement chapter~\ref{ch:data-without-n-plus-1}'s interceptor prints. What
makes it survivable is that it reaches the client, which means it reaches
whatever the client reports, and an error a caller can see is an error somebody
eventually raises. The wrong-order half does not even do that. It will sit in a
graph until a person notices that a talk has the wrong speaker's name on it and
is believed.

\subsection{What the Defect Did Not Reach}

\index{DataLoader!de-duplication is above the defect}%
Two properties of the entity route survive this loader untouched, and they are
part of why the file reads correctly.

De-duplication still works. Send the same representation three times and one
key goes into the \code{WHERE} clause, and three correct copies of the same
speaker come back. That is because the de-duplication lives in
\code{DataLoaderBase} above \code{FetchAsync}, in the part nobody replaced, and
a batch of one identical key cannot be out of order.

\index{reference resolver!the guard is above the defect too}%
The type-name guard still works. Send a \code{Session} node id under a
\code{Speaker} type name in the middle of two good keys and the middle entry is
null, two keys reach the batch, and the neighbors are untouched. That guard is
in the reference resolver, which is the file it always was.

\index{representation order!narrow is not the same as safe}%
So the defect is narrow. It touches one method, it leaves every other property
of the entity route intact, and the file it lives in reads correctly. Narrow is
not the same as safe: this is the smallest change I
know of that makes a federated graph lie to a client and keep every one of its
instruments green.
\index{representation order!the ascending coincidence|)}%
